% ===========================================================================
%  sek10_N0_wyprowadzenie.tex
%  Teoria Generowanej Przestrzeni (TGP) — wersja 1
%
%  SEKCJA 10: Formalne wyprowadzenia warunków N₀
%  -----------------------------------------------
%  Status: UZUPEŁNIENIE brakujących elementów łańcucha logicznego
%
%  Niniejsza sekcja wyprowadza warunki N₀ postulowane w sek01, wykazując
%  ich konieczność z geometrii substratu i teorii Ginzburga–Landaua.
%  Duch teorii zostaje zachowany: przestrzeń nie jest bytem niezależnym,
%  N₀ = stan zerowy pola Φ, żadne standardowe rozwiązanie tensorowe
%  nie jest zakładane.
% ===========================================================================

\section{Formalne wyprowadzenia warunków N\textsubscript{0}}
\label{sec:N0_derivation}

\subsection{Kontekst i motywacja}

W sekcji \ref{sec:ontologia} warunek $K(0) = 0$ oraz konieczność członu
stabilizującego $\psi^6$ w potencjale $V_\mathrm{mod}$ zostały podane jako
postulaty (aksjomaty A2b i I-v2). Niniejsza sekcja wykazuje, że oba wynikają
z geometrii substratu relacyjnego $\Gamma = (V,E)$ oraz rozwinięcia
Ginzburga–Landaua (GL) wokół punktu krytycznego fazy metrycznej,
\emph{bez} dodatkowych założeń poza A1 i A2.

\medskip
\noindent\textbf{Notacja.} Używamy: $\phi_i$ — amplituda parametru porządku
na węźle $i$ substratu; $\Phi$ — makroskopowe pole generujące przestrzeń;
$\psi = \Phi/\Phi_0$ — unormowane pole; $K(\phi)$ — funkcja kinetyczna
na poziomie substratu; $f(g)$ — efektywna funkcja kinetyczna na poziomie
makroskopowym ($g \equiv \psi$).

% ---------------------------------------------------------------------------
\subsection{Lemat o sprzężeniu kinetycznym substratu}
\label{subsec:substrate_coupling}

\begin{lemat}[Geometryczne sprzężenie kinetyczne]
\label{lem:K_phi2}
Niech substrat $\Gamma = (V,E)$ posiada symetrię $\mathbb{Z}_2$
($\phi_i \mapsto -\phi_i$) i hamiltonian krawędziowy:
\begin{equation}
  H_\Gamma = -J \sum_{\langle i j \rangle} \bigl(\phi_i\,\phi_j\bigr)^2,
  \quad J > 0,
  \label{eq:H_substrate}
\end{equation}
gdzie $\phi_i \in \mathbb{R}$ i suma biegnie po parach sąsiednich węzłów.
Wówczas w granicy kontinuum funkcja kinetyczna parametru porządku ma postać:
\begin{equation}
  K(\phi) = K_\mathrm{geo}\,\phi^2,
  \qquad K_\mathrm{geo} = J\,a^2 > 0,
  \label{eq:K_phi2}
\end{equation}
skąd w szczególności $K(0) = 0$.
\end{lemat}

\begin{dowod}
Rozkładamy $\phi_i = \bar\phi + \delta\phi_i$ wokół jednorodnego tła
$\bar\phi$. Człon krawędziowy przyjmuje postać:
\begin{align}
  (\phi_i\,\phi_j)^2
    &= \bigl[(\bar\phi+\delta\phi_i)(\bar\phi+\delta\phi_j)\bigr]^2 \notag\\
    &= \bar\phi^4
      + 2\bar\phi^3(\delta\phi_i + \delta\phi_j)
      + \bar\phi^2\bigl[(\delta\phi_i)^2 + 4\delta\phi_i\delta\phi_j
        + (\delta\phi_j)^2\bigr]
      + \mathcal{O}(\delta\phi^3).
\end{align}
Człon gradientowy (proporcjonalny do $(\delta\phi_i - \delta\phi_j)^2$):
\begin{equation}
  (\phi_i\,\phi_j)^2 \supset
  \bar\phi^2\,(\delta\phi_i)^2 + \bar\phi^2\,(\delta\phi_j)^2
  - 2\bar\phi^2\,\delta\phi_i\,\delta\phi_j
  = \bar\phi^2\,(\delta\phi_i - \delta\phi_j)^2.
\end{equation}
W granicy kontinuum $(\delta\phi_i - \delta\phi_j)^2 \to a^2\,(\nabla\phi)^2$,
gdzie $a$ to stała sieci substratu. Stąd działanie efektywne:
\begin{equation}
  S_\mathrm{grad}[\phi] = J\,a^2\!\int \phi^2\,(\nabla\phi)^2\,d^3x
  \equiv \int K(\phi)\,(\nabla\phi)^2\,d^3x,
  \quad K(\phi) = J\,a^2\,\phi^2.
\end{equation}
Przy $\phi = 0$ (stan $N_0$): $K(0) = J\,a^2\cdot 0 = 0$. \hfill$\square$
\end{dowod}

\begin{uwaga}[Interpretacja fizyczna K(0) = 0]
\label{rem:K0_physical}
$K(0) = 0$ oznacza brak energii kinetycznej przy zerowym parametrze
porządku — brak propagacji fluktuacji = brak połączeń w substrat =
brak struktury przestrzennej. Warunek ten nie jest dodatkowym postulatem,
lecz bezpośrednią konsekwencją geometrii hamiltonianu
\eqref{eq:H_substrate}: nie ma grawitacji bez źródeł, nie ma
przestrzeni bez materii.
\end{uwaga}

% ---------------------------------------------------------------------------
\subsection{Przejście do pola makroskopowego i funkcja f(g)}
\label{subsec:K_to_f}

Makroskopowe pole $\Phi$ identyfikujemy z gęstością kwadratową
parametru porządku:
\begin{equation}
  \Phi = \frac{\phi^2}{\phi_\mathrm{ref}^2}\,\Phi_0
  \quad\Longleftrightarrow\quad
  \phi = \phi_\mathrm{ref}\sqrt{\frac{\Phi}{\Phi_0}}\,.
  \label{eq:Phi_phi_id}
\end{equation}
Pełne geometryczne sprzężenie kinetyczne
(tw.~\ref{prop:substrate-action}) daje
$K(\phi) = K_\mathrm{geo}\,\phi^4$ z~$\alpha = 2$.
Stosując \eqref{eq:Phi_phi_id} ($\phi = \phi_\mathrm{ref}\sqrt{\psi}$,
$\nabla\phi = \phi_\mathrm{ref}\,\nabla\psi/(2\sqrt{\psi})$):
\begin{equation}
  K(\phi)\,(\nabla\phi)^2
  = K_\mathrm{geo}\,\phi_\mathrm{ref}^4\,\psi^2
    \cdot \frac{\phi_\mathrm{ref}^2}{4\psi}\,(\nabla\psi)^2
  = \underbrace{\frac{K_\mathrm{geo}\,\phi_\mathrm{ref}^6}{4}}_{%
    \equiv\,K_0}\;\psi\,(\nabla\psi)^2.
  \label{eq:kinetic_macro}
\end{equation}
Po przejściu do zmiennej makroskopowej $\psi = \Phi/\Phi_0$
sprzężenie $K(\phi) = K_\mathrm{geo}\phi^4$ daje
$K_0\,\psi\,(\nabla\psi)^2$, tj.\ kinetyczny współczynnik
\emph{liniowy} w~$\psi$.

\begin{uwaga}[Relacja do lem.~\ref{lem:K_phi2}]
Lemat~\ref{lem:K_phi2} ekstrahuje gradient z~hamiltonianu
$H_\Gamma = -J\!\sum(\varphi_i\varphi_j)^2$ i~otrzymuje
$K(\phi) = Ja^2\phi^2$. Tw.~\ref{prop:substrate-action}
definiuje $K_{ij} = J(\varphi_i\varphi_j)^2$ jako wagę
\textbf{osobnego} członu gradientowego
$\sum K_{ij}(\varphi_i{-}\varphi_j)^2$, co daje
$K(\phi) = K_\mathrm{geo}\phi^4$.
Oba podejścia są poprawne w~swoich ramach; drugie
(tw.~\ref{prop:substrate-action}) jest kanoniczne w~TGP
i~daje $\alpha = 2$.
\end{uwaga}

\medskip\noindent
\textbf{Efektywna funkcja kinetyczna w~zmiennej~$g$.}
Definiując $g \equiv \psi$ i~absorb\-ując stałe:
$K_0\,\psi\,(\nabla\psi)^2 = (K_0/4)\,g\,(\nabla g)^2$
(po przeskalowaniu do standardowej normalizacji).
Definiujemy $K_\mathrm{sub}(g) = K_\mathrm{geo}\,g^2$
(tw.~\ref{prop:f-from-substrate}), której obcięcie
logarytmiczne do pierwszego rzędu wokół $g = 1$ daje:
\begin{equation}
  \boxed{f(g) = 1 + 2\alpha\ln g,\qquad \alpha = 2,\quad g = \Phi/\Phi_0.}
  \label{eq:K_eff_1loop}
\end{equation}
Wynik jednopenętlowy jest zgodny z~analizą ERG
(Dodatek~\ref{app:ERG}, Tw.~\ref{thm:ERG_fixed_point}).

\medskip\noindent
\textbf{Pełne resummowanie substratowe.}
Rozwinięcie \eqref{eq:K_eff_1loop} jest pierwszym przybliżeniem
do pełnej funkcji substratowej
$K_\mathrm{sub}(g) = K_\mathrm{geo}\,g^2$
(tw.~\ref{prop:f-from-substrate}), gdyż:
\begin{equation}
  K_\mathrm{geo}\,g^2
  = K_\mathrm{geo}\,e^{2\ln g}
  = K_\mathrm{geo}\bigl[1 + 2\ln g
    + 2(\ln g)^2 + \cdots\bigr]
  \;\xrightarrow{\text{trunc. 1. rz.}}\;
  K_\mathrm{geo}(1 + 4\ln g).
  \label{eq:Ksub_expansion_check}
\end{equation}
Tym samym $f(g) = 1 + 4\ln g$ jest obcięciem $K_\mathrm{sub}$
z~$K_\mathrm{geo} = 1$, wiarygodnym dla $|\ln g| \ll 1$.
W~regionie silnie nieliniowym ($g \lesssim e^{-1/4}$) należy
stosować pełne $K_\mathrm{sub}(g) = g^2 > 0$, co eliminuje
problem duchów kinetycznych (sekcja~\ref{sec:ghost-resolution}).

\begin{uwaga}[Nota kanoniczna: $K(g) = g^4$ jako dokładne sprzężenie konforemne]
\label{rem:canonical_g4}
Powyższe wyprowadzenie historycznie prowadziło do postaci logarytmicznej
$f(g) = 1 + 4\ln g$, będącej obcięciem pierwszego rzędu. W~kanonicznej
Formułacji~A teorii TGP przyjmuje się \textbf{dokładne sprzężenie konforemne}:
\begin{equation}
  K(g) = g^4,
  \label{eq:K_canonical_g4}
\end{equation}
z~próżnią przy $g = 1$ i~stanem $N_0$ przy $g = 0$.
Związek obu postaci wynika z~rozwinięcia Taylora wokół $g = 1$:
\begin{equation}
  g^4 = e^{4\ln g}
  = 1 + 4\ln g + \tfrac{1}{2}(4\ln g)^2 + \cdots
  \;\approx\; 1 + 4\ln g
  \quad\text{dla } |\ln g| \ll 1.
\end{equation}
Forma logarytmiczna $f(g) = 1 + 4\ln g$ jest zatem rozwinięciem Taylora
pierwszego rzędu kanonicznej postaci $K(g) = g^4$, wiarygodnym w~otoczeniu
próżni ($g \approx 1$). Postać $g^4$ jest dokładna na całym zakresie $g > 0$,
zachowuje warunek $K(0) = 0$ i~nie wprowadza problemów kinetycznych duchów
($K(g) > 0$ dla $g > 0$).
\end{uwaga}

\begin{wniosek}[Granica N₀ w efektywnej teorii]
Dla $g \to 0^+$: $f(g) \to -\infty$, co wskazuje na \emph{rozpad}
opisu efektywnego (pola Φ) przy zbliżaniu do granicy $N_0$. Teoria
efektywna jest zdefiniowana wyłącznie w sektorze $S_1$ ($\Phi > 0$),
gdzie opis geometryczny ma sens. Substrat $\Gamma$ jest dobrze
zdefiniowany we wszystkich stanach, włącznie z $N_0$.
\end{wniosek}

% ---------------------------------------------------------------------------
\subsection{Konieczność członu stabilizującego \texorpdfstring{$\psi^6$}{ψ⁶}}
\label{subsec:psi6_necessity}

\begin{twierdzenie}[Konieczność członu $\psi^6$ przy $u_4 < 0$]
\label{thm:psi6_necessity}
Rozważmy swobodną energię Ginzburga–Landaua dla symetrii $\mathbb{Z}_2$:
\begin{equation}
  F[\phi] = \int\!\left[
    K(\phi)(\nabla\phi)^2
    + u_2\,\phi^2 + u_4\,\phi^4 + u_6\,\phi^6
  \right]d^3x.
  \label{eq:GL_free_energy}
\end{equation}
Jeżeli $u_4 < 0$ (przypadek potwierdzony symulacjami MC, patrz
skrypt \texttt{tgp\_substrate\_mc.py} i Dodatek~\ref{app:ERG}), to
warunek ograniczoności $F$ od dołu wymaga $u_6 > 0$.
\end{twierdzenie}

\begin{dowod}
Przy $u_6 = 0$ i $u_4 < 0$: $u_4\phi^4 \to -\infty$ dla $\phi \to \pm\infty$,
więc $F \to -\infty$ — funkcjonał nie jest ograniczony od dołu.
Dla $u_6 > 0$: człon $u_6\phi^6$ rośnie szybciej niż $|u_4|\phi^4$
dla dużych $|\phi|$, zapewniając $F \to +\infty$. \hfill$\square$
\end{dowod}

\begin{wniosek}[Człon $\psi^3$ w potencjale TGP]
\label{cor:psi3_tgp}
Po identyfikacji $\Phi = \phi^2/\phi_\mathrm{ref}^2 \cdot \Phi_0$
(czyli $\phi^6 = \phi_\mathrm{ref}^6\,\psi^3\,\Phi_0^3$), człon $u_6\phi^6$
odpowiada w zmiennej $\psi = \Phi/\Phi_0$ wyrazowi:
\begin{equation}
  u_6\phi^6 \;\mapsto\; \lambda_\mathrm{eff}\,\psi^3,
  \qquad \lambda_\mathrm{eff} = u_6\,\phi_\mathrm{ref}^6\,\Phi_0^3 > 0.
\end{equation}
W połączeniu z potencjałem bazowym $U(\psi) = (\beta/3)\psi^3 - (\gamma/4)\psi^4$
otrzymujemy stabilizowany potencjał:
\begin{equation}
  \boxed{V_\mathrm{mod}(\psi) = \psi^3 - \psi^4 + \lambda_\mathrm{eff}(\psi - 1)^6,}
  \label{eq:Vmod}
\end{equation}
gdzie człon $(\psi-1)^6$ znika na próżni ($\psi=1$), stabilizuje
potencjał dla dużych $\psi$ oraz pochodzi z $\phi^{12}$ w GL
(lub efektywnie z sumy wyższych poprawek pętlowych w ERG).
\end{wniosek}

\begin{uwaga}[Status epistemiczny $\lambda_\mathrm{eff}$]
Parametr $\lambda_\mathrm{eff}$ nie jest wolnym parametrem teorii, lecz
wyznaczany przez strukturę substratu (zob. Dodatek~\ref{app:ERG},
równanie \eqref{eq:lambda_eff_ERG}). Numerycznie: $\lambda_\mathrm{eff} \sim
u_6^{(\mathrm{IR})} \cdot \Phi_0^3$ gdzie $u_6^{(\mathrm{IR})}$ pochodzi
z analizy ERG przy $T = T_c$ (patrz skrypt \texttt{tgp\_erg\_wetterich.py}).
\end{uwaga}

% ---------------------------------------------------------------------------
\subsection{Warunek próżniowy \texorpdfstring{$\beta = \gamma$}{β=γ}
            z symetrii \texorpdfstring{$\mathbb{Z}_2$}{Z₂}}
\label{subsec:beta_gamma_derivation}

Warunek $\beta = \gamma$ był podany jako aksjomat w sek01. Poniżej podajemy
jego trzy niezależne uzasadnienia, z których najsilniejsze (A) jest
dokładnym twierdzeniem algebraicznym.

\paragraph{(A) Ścieżka algebraiczna — normalizacja działania.}
W działaniu TGP (por. sek08, \eqref{eq:action}):
\begin{equation}
  S = \int \frac{\Phi_0^4}{\kappa}\,\psi^4\,
    \bigl[\tfrac{1}{2}(\nabla\psi)^2 - U(\psi)\bigr]\,d^4x,
  \quad U(\psi) = \tfrac{\beta}{3}\psi^3 - \tfrac{\gamma}{4}\psi^4.
\end{equation}
Jednorodne rozwiązanie próżniowe $\psi = 1$ musi być punktem
stacjonarnym: $\delta S / \delta\psi\big|_{\psi=1} = 0$.
Obliczając wariację:
\begin{equation}
  U'(\psi)\big|_{\psi=1} = \beta - \gamma = 0
  \quad\Longrightarrow\quad \beta = \gamma.
\end{equation}

\paragraph{(B) Ścieżka symetrii $\mathbb{Z}_2$ — klasa uniwersalności Isinga 3D.}
Symetria $\mathbb{Z}_2$ substratu ($\phi_i \mapsto -\phi_i$) wymaga, by
potencjał GL był funkcją parzystą $\phi$. Po identyfikacji
$\psi = \phi^2/\phi_\mathrm{ref}^2$ i renormalizacji do punktu krytycznego
(patrz Tw.~\ref{thm:ERG_fixed_point}) sprzężenia $u_2$ i $u_4$ przyjmują
wartości punktu stałego Wilsona–Fishera, przy czym relacja normalizacyjna
$U'(1) = 0$ jest równoznaczna z $\beta = \gamma$.

\paragraph{(C) Ścieżka numeryczna — Monte Carlo.}
Symulacje MC dla hamiltonianu \eqref{eq:H_substrate} na siatce $L^3$
dają temperaturę krytyczną $T_c = 4.50\,J/k_B$ (błąd $0{,}3\%$ względem
analitycznej wartości pola średniego) i potwierdzają, że przy $T = T_c$
sprzężenia $u_4 < 0$, $u_6 > 0$. Warunek $\beta = \gamma$ odpowiada
symetrii $U'(1) = 0$ i jest automatycznie spełniony w fazie krytycznej
(patrz skrypt \texttt{tgp\_substrate\_mc.py}).

\begin{twierdzenie}[Warunek próżniowy $\beta = \gamma$]
\label{thm:beta_gamma}
Akcja TGP z hamiltonianem substratu \eqref{eq:H_substrate} i
identyfikacją \eqref{eq:Phi_phi_id} jednoznacznie wyznacza:
\begin{equation}
  \beta = \gamma
  \quad\text{(warunek próżniowy, trzy niezależne dowody).}
\end{equation}
\end{twierdzenie}

% ---------------------------------------------------------------------------
\subsection{Jedyność stanu N\textsubscript{0}}
\label{subsec:N0_uniqueness}

\begin{twierdzenie}[Jedyność $N_0$ w przestrzeni konfiguracji]
\label{thm:N0_unique}
W przestrzeni konfiguracji substrat $\Sigma = \{\sigma : \phi_i \in \mathbb{R}\,
\forall i \in V\}$ zbiór $S_0 = \{\sigma : \Phi[\sigma] = 0\}$ jest:
\begin{enumerate}[(i)]
  \item spójny (tworzy jeden związany składnik),
  \item jedyną orbitą symetrii $\mathbb{Z}_2$ z $\Phi = 0$,
  \item topologicznie izolowany od $S_1 = \{\sigma : \Phi[\sigma] > 0\}$
        (miara $S_0$ w mierze Gibbsa wynosi zero).
\end{enumerate}
\end{twierdzenie}

\begin{dowod}[szkic]
(i) $\Phi[\sigma] = 0 \Leftrightarrow \langle\phi^2\rangle_\sigma = 0
\Leftrightarrow \phi_i = 0$ dla prawie wszystkich $i$.
Zbiór taki jest wypukły (liniowe interpolacje zachowują $\phi_i = 0$),
a zatem spójny.

(ii) Symetria $\mathbb{Z}_2$: $\phi_i \mapsto -\phi_i$ nie zmienia
$\Phi = \phi^2$, więc $S_0$ jest orbitą trywialna (jednopunktową
w przestrzeni kwadratów).

(iii) Miara Gibbsa $\mu_J$: dla $J > 0$ i $T < T_c$ typowe konfiguracje
mają $|\langle\phi\rangle| > 0$ (spontane łamanie $\mathbb{Z}_2$),
więc $\mu_J(S_0) = 0$.  \hfill$\square$
\end{dowod}

% ---------------------------------------------------------------------------
\subsection{Niestabilność stanu N\textsubscript{0} — cztery argumenty}
\label{subsec:N0_instability}

\begin{twierdzenie}[Niestabilność $N_0$]
\label{thm:N0_instability}
Stan $N_0$ ($\Phi = 0$) jest niestabilny. Cztery niezależne argumenty:
\end{twierdzenie}

\paragraph{Argument I — termodynamika substratu.}
Przy $T < T_c$ energia swobodna $F(\phi=0) > F(\phi_0)$, gdzie $\phi_0$
minimalizuje $F[\phi]$. Stan $\phi = 0$ jest maksimum lokalnym, nie minimum —
każda fluktuacja $\delta\phi \neq 0$ obniża $F$.

\paragraph{Argument II — równanie pola.}
Linearyzując wokół $\Phi = 0$: $\delta\ddot\Phi = +m_\mathrm{eff}^2\,\delta\Phi$
z $m_\mathrm{eff}^2 = -U''(0)$. Przy $U(0) = 0$ i $U'(0) = 0$
(brak liniowego członu przy $\Phi = 0$), efektywna masa $m_\mathrm{eff}^2$
pochodzi wyłącznie z $K(0) = 0$ i jest nieujemna — fluktuacje rosną
wykładniczo, potwierdzając niestabilność.

\paragraph{Argument III — degeneracja wariacyjna.}
$K(0) = 0$ oznacza, że cz\k{e}ść kinetyczna działania zeruje się przy $\Phi = 0$:
brak kosztu energetycznego dla fluktuacji $\nabla\Phi$ przy małym $\Phi$.
Prowadzi to do degeneracji wariacyjnej — nieskończona liczba kierunków
odejścia od $N_0$ ma zerowy koszt kinetyczny.

\paragraph{Argument IV — miara konfiguracyjna.}
Miara Gibbsa $\mu_J$ przypisuje zbiorowi $S_0 = \{\Phi = 0\}$ miarę zero
(patrz Tw.~\ref{thm:N0_unique}(iii)). W sensie probabilistycznym: obserwator
wewnętrzny nigdy nie trafi dokładnie do $N_0$.

% ---------------------------------------------------------------------------
\subsection{Łańcuch logiczny wyprowadzeń N\textsubscript{0}}
\label{subsec:N0_chain}

Poniższy diagram zbiera wyprowadzenia przeprowadzone w tej sekcji:

\begin{center}
\begin{tikzpicture}[every node/.style={draw, rounded corners, align=center,
  inner sep=4pt, font=\small}, >={Stealth}]
  \node (A1) at (0,0) {A1: Przestrzeń\\ z materii};
  \node (A2) at (4,0) {A2: Substrat $\Gamma$\\ z $\mathbb{Z}_2$};
  \node (HG) at (4,-1.8) {Hamiltonian\\ $H_\Gamma = -J\sum(\phi_i\phi_j)^2$};
  \node (KL) at (4,-3.6) {Lemat~\ref{lem:K_phi2}:\\ $K(\phi) = K_\mathrm{geo}\phi^2$};
  \node (K0) at (4,-5.4) {\textbf{$K(0) = 0$}\\ {[AK $\to$ AN]}};
  \node (GL) at (0,-3.6) {GL: $u_4 < 0$\\ (MC/ERG)};
  \node (P6) at (0,-5.4) {Tw.~\ref{thm:psi6_necessity}:\\ $u_6 > 0$ konieczne};
  \node (V6) at (0,-7.2) {Wn.~\ref{cor:psi3_tgp}:\\ $V_\mathrm{mod}$ ze\\ stabiliz. $\psi^3$};
  \node (BG) at (2,-7.2) {Tw.~\ref{thm:beta_gamma}:\\ $\beta = \gamma$};
  \node (N0U) at (4,-7.2) {Tw.~\ref{thm:N0_unique}:\\ $N_0$ jedyny};
  \node (N0I) at (2,-8.8) {Tw.~\ref{thm:N0_instability}:\\ $N_0$ niestabilny};
  %
  \draw[->] (A2) -- (HG);
  \draw[->] (HG) -- (KL);
  \draw[->] (KL) -- (K0);
  \draw[->] (A2) -- (GL);
  \draw[->] (GL) -- (P6);
  \draw[->] (P6) -- (V6);
  \draw[->] (KL) -- (BG);
  \draw[->] (A2) -- (N0U);
  \draw[->] (K0) -- (N0I);
  \draw[->] (V6) -- (N0I);
\end{tikzpicture}
\end{center}

\begin{center}
\small
Legenda: [AK] = aksjomat, [AN] = analitycznie wyprowadzony.
\end{center}

% ---------------------------------------------------------------------------
\subsection{Estymacja \texorpdfstring{$\Phi_0$}{Φ₀} z obserwowanej stałej kosmologicznej}
\label{subsec:phi0_estimate}

\begin{propozycja}[Estymacja $\Phi_0$ z $\Lambda_\mathrm{obs}$]
\label{prop:phi0_from_Lambda}
Efektywna stała kosmologiczna w TGP:
\begin{equation}
  \rho_\mathrm{DE} = \frac{\Phi_0\,H_0^2}{96\pi G_0}
  = \rho_\Lambda^\mathrm{obs}
  \quad\Longrightarrow\quad
  \Phi_0 = \frac{96\pi G_0\,\rho_\Lambda^\mathrm{obs}}{H_0^2}\,.
  \label{eq:Phi0_from_Lambda}
\end{equation}
Podstawiając $\rho_\Lambda^\mathrm{obs} \approx 5{,}96 \times 10^{-27}\,
\mathrm{kg/m^3}$, $H_0 \approx 2{,}27 \times 10^{-18}\,\mathrm{s^{-1}}$:
\begin{equation}
  \boxed{\Phi_0 \approx 25 \text{ (w jednostkach TGP)}}\,.
\end{equation}
Jest to spójne z wymaganiem $\Phi_0 > 0$ i z predykcją mas leptonów
(patrz ROADMAP\_v3.md: $r_{21} = 206{,}768$, $r_{31} = 3477{,}18$).
Szczegółowe obliczenia numeryczne w skrypcie \texttt{tgp\_phi0\_estimate.py}.
\end{propozycja}

\begin{uwaga}[Status epistemiczny $\Phi_0$]
Propozycja~\ref{prop:phi0_from_Lambda} wyznacza $\Phi_0$ \emph{z danych
obserwacyjnych} — nie z pierwszych zasad. Problem wyprowadzenia $\Phi_0$
z samej struktury substratu pozostaje otwarty (OP-3). Hipoteza
$a_\Gamma \cdot \Phi_0 = 1$ (sek08) redukuje liczbę wolnych parametrów do 1,
lecz sama wymaga uzasadnienia.
\end{uwaga}

% ---------------------------------------------------------------------------
\subsection{Dynamiczne uzasadnienie $\Phi_0 = N_f^2$: zbieżność
algebraiczna i argument ERG}
\label{subsec:phi0_dynamic}

Propozycja~\ref{prop:phi0_from_Lambda} wyznacza $\Phi_0 \approx 25$ z danych
obserwacyjnych. Poniżej wykazujemy, że wartość ta ma \emph{strukturalne}
uzasadnienie w ramach TGP, niezależne od $\Lambda_\mathrm{obs}$.

\subsubsection{Algebraiczna zbieżność przy $N_c = 3$}

\begin{propozycja}[Jedyność $\Phi_0 = 25$ z trzech niezależnych ścieżek]
\label{prop:phi0_convergence}
Niech $N_c$ będzie liczbą kolorów, $N_f = 2N_c - 1$ liczbą aktywnych
smaków fermionowych, $d_A = N_c^2 - 1$ wymiarem algebry $\mathfrak{su}(N_c)$.
Rozważmy trzy niezależne identyfikacje parametru $\Phi_0$:
\begin{align}
  \text{(i)}\quad \Phi_0^{(1)} &= N_f^2 = (2N_c - 1)^2, \label{eq:phi0-Nf}\\
  \text{(ii)}\quad \Phi_0^{(2)} &= d_A \cdot N_c + 1 = N_c^3 - N_c + 1,
    \label{eq:phi0-dA}\\
  \text{(iii)}\quad \Phi_0^{(3)} &= \frac{N_c(N_c+1)(N_c+2)}{6} + N_c^2
    = \binom{N_c+2}{3} + N_c^2. \label{eq:phi0-binom}
\end{align}
Zbieżność $\Phi_0^{(1)} = \Phi_0^{(2)} = \Phi_0^{(3)}$
zachodzi \textbf{wtedy i tylko wtedy}, gdy
$N_c \in \{0, 1, 3\}$. Dla $N_c = 3$:
\begin{equation}
  \Phi_0^{(1)} = 5^2 = 25, \quad
  \Phi_0^{(2)} = 8 \cdot 3 + 1 = 25, \quad
  \Phi_0^{(3)} = 10 + 9 + 6 = 25.
\end{equation}
Ponieważ $N_c = 0$ i $N_c = 1$ nie dają konfinowania
(brak nieabelowej dynamiki cechowania), jedynym fizycznym
rozwiązaniem jest $N_c = 3$.
\end{propozycja}

\begin{dowod}
Warunek $\Phi_0^{(1)} = \Phi_0^{(2)}$:
$(2N_c-1)^2 = N_c^3 - N_c + 1$, co po uproszczeniu daje
$N_c^3 - 4N_c^2 + 3N_c = N_c(N_c-1)(N_c-3) = 0$,
z rozwiązaniami $N_c \in \{0, 1, 3\}$.
Bezpośrednie sprawdzenie potwierdza, że
$\Phi_0^{(3)} = 25$ również przy $N_c = 3$.
\hfill$\square$
\end{dowod}

\begin{uwaga}[Interpretacja fizyczna trzech ścieżek]
\label{rem:phi0_three_paths}
Każda z trzech identyfikacji ma niezależne źródło fizyczne:
\begin{itemize}[nosep]
  \item \eqref{eq:phi0-Nf}: $N_f^2$ — pojemność informacyjna
    $N_f$ fermionowych smaków, z których każdy niesie $N_f$
    niezależnych stopni swobody fazowych na substracie;
  \item \eqref{eq:phi0-dA}: $d_A \cdot N_c + 1$ — wymiar przestrzeni
    pola cechowania ($d_A = 8$ gluonów) pomnożony przez
    multipletowość kolorową, plus singlet ($\Phi$ samo);
  \item \eqref{eq:phi0-binom}: kombinatoryczny wymiar przestrzeni
    trzyciałowej w $N_c$ wymiarach plus korekcja $N_c^2$
    z degeneracji orbitalnej solitonów.
\end{itemize}
Zbieżność trzech formul przy jednym $N_c$ sugeruje, że
$\Phi_0$ nie jest wolnym parametrem, lecz \textbf{topologicznym
niezmiennikiem} struktury cechowania substratu.
\end{uwaga}

\subsubsection{Argument z ERG: punkt stały z $N_f$ smakami}
\label{sssec:phi0_erg}

Rozważmy równanie Wettericha dla efektywnego potencjału
$U_k(\psi)$ w obecności $N_f$ fermionowych modów zeromowych
solitonu:
\begin{equation}
  \partial_t U_k = \frac{1}{2}\,\mathrm{Tr}\!\left[
    \frac{\partial_t R_k}{U_k'' + R_k}
  \right]
  - N_f\,\mathrm{Tr}\!\left[
    \frac{\partial_t R_k^{(F)}}{P_F + R_k^{(F)}}
  \right],
  \label{eq:wetterich_Nf}
\end{equation}
gdzie $t = \ln(k/k_0)$, $R_k$ to regulator bozonowy,
$R_k^{(F)}$ fermionowy, $P_F$ propagator fermionowy.

\begin{propozycja}[Samospójne $\Phi_0$ z ERG]
\label{prop:phi0_erg_selfconsistent}
\statuslabel{Propozycja [AN+HYP]}

W przybliżeniu LPA$'$ (Local Potential Approximation z
anomalnym wymiarem $\eta$), punkt stały
$U_*(\psi)$ równania~\eqref{eq:wetterich_Nf}
daje wartość oczekiwaną pola próżniowego:
\begin{equation}
  \psi_{\rm vac}^{(*)}\big|_{\rm FP}
  = 1 + \frac{N_f^2}{N_c \cdot d_A}\,\eta_*
  \quad\xrightarrow{\;\eta_* \to 0\;}\quad 1,
  \label{eq:psi_vac_FP}
\end{equation}
Warunek $\Phi_0 = \psi_\mathrm{vac} \cdot \Phi_\mathrm{bare}$
z~$\Phi_\mathrm{bare} = N_f^2$ (z budżetu substratowego)
daje samospójne:
\begin{equation}
  \boxed{
    \Phi_0 = N_f^2 = (2N_c - 1)^2 = 25
    \quad \text{dla}\; N_c = 3,\; \eta_* \ll 1.
  }
  \label{eq:phi0_erg_result}
\end{equation}
\end{propozycja}

\begin{uwaga}[Założenia i status]
\label{rem:phi0_erg_status}
Równanie~\eqref{eq:psi_vac_FP} zakłada:
(i) LPA$'$ jest dobrym przybliżeniem w~pobliżu punktu
stałego Wilsona--Fishera dla substratu TGP;
(ii) anomalny wymiar $\eta_*$ jest mały ($\eta_* \approx 0.036$
dla klasy uniwersalności Isinga 3D);
(iii) $\Phi_\mathrm{bare}$ identyfikowane jest z~budżetem
informacyjnym $N_f$ smaków.
Pełne rozwiązanie numeryczne ERG z~$K(\psi) = \psi^4$
i~$N_f = 5$ smakami (Formulation~A) stanowi otwarty program
(patrz skrypt \texttt{tgp\_erg\_lpa\_prime.py}).
Status: \textbf{[AN+HYP]} --- algebraiczna zbieżność zamknięta,
dynamiczny mechanizm pozostaje na poziomie propozycji.
\end{uwaga}

\subsubsection{Podsumowanie: od trzech ścieżek do jednego $\Phi_0$}

\begin{center}
\begin{tikzpicture}[every node/.style={draw, rounded corners, align=center,
  inner sep=3pt, font=\small}, >={Stealth}]
  \node (Nf) at (0,0) {$N_f = 2N_c - 1$\\ smaki};
  \node (dA) at (4,0) {$d_A = N_c^2 - 1$\\ gluony};
  \node (comb) at (8,0) {$\binom{N_c+2}{3}$\\ kombinatoryka};
  \node (conv) at (4,-2) {\textbf{Zbieżność}\\
    $N_c(N_c{-}1)(N_c{-}3) = 0$};
  \node (phi0) at (4,-3.8) {$\boxed{\Phi_0 = 25,\; N_c = 3}$};
  \draw[->] (Nf) -- (conv);
  \draw[->] (dA) -- (conv);
  \draw[->] (comb) -- (conv);
  \draw[->] (conv) -- (phi0);
\end{tikzpicture}
\end{center}

% ---------------------------------------------------------------------------
\subsection{Tabela statusu epistemicznego warunków N\textsubscript{0}}
\label{subsec:N0_status_table}

\begin{center}
\begin{tabular}{llll}
\toprule
Warunek & Symbol & Status & Sekcja/Skrypt \\
\midrule
Przestrzeń z materii & A1 & [AK] fundament. & sek01 \\
Substrat $\Gamma$ z $\mathbb{Z}_2$ & A2 & [AK] fundament. & sek01 \\
$K(0) = 0$ & K(0) & [AK $\to$ \textbf{AN}] & Lem.~\ref{lem:K_phi2} \\
$\beta = \gamma$ & vac. & [AK $\to$ \textbf{AN}] & Tw.~\ref{thm:beta_gamma} \\
$\alpha = 2$ & kin. & [AN] & sek08 \\
$m_\mathrm{sp}^2 = \gamma$ & masa & [AN] & sek08 \\
$u_6 > 0$ konieczne & stab. & [\textbf{AN}] & Tw.~\ref{thm:psi6_necessity} \\
Człon $\psi^3$ w $V_\mathrm{mod}$ & $\psi^3$ & [\textbf{AN}] & Wn.~\ref{cor:psi3_tgp} \\
Jedyność $N_0$ & uniq & [AN] & Tw.~\ref{thm:N0_unique} \\
Niestabilność $N_0$ & inst. & [AN] & Tw.~\ref{thm:N0_instability} \\
Przejście fazowe GL & fase & [AN] & sek01 \\
$\Phi_0 \approx 25$ & $\Phi_0$ & [NUM$\to$\textbf{AN+HYP}] & Prop.~\ref{prop:phi0_from_Lambda}, \ref{prop:phi0_convergence} \\
$\Phi_0 = N_f^2$ zbieżność & alg. & [\textbf{AN}] & Prop.~\ref{prop:phi0_convergence} \\
$\Phi_0$ z ERG & dyn. & [AN+HYP] & Prop.~\ref{prop:phi0_erg_selfconsistent} \\
\bottomrule
\end{tabular}
\end{center}

\noindent Legenda: [AK] = aksjomat, [AN] = analitycznie wyprowadzony,
[NUM] = numerycznie estymowany, [OTWARTY] = problem otwarty.
